\section{Affine minima with arbitrary tolerance}
\label{sec:affine}

The proof will use a finite minimum of affine functions on the positive
orthant. Its essential property is that an active slope approximates the
exponent vector of the point where it is active. The allowed error will
depend on the slope and may be discontinuous. We prove the construction
and its uniform consequence from
\cite[Lemma~2.1 and Corollary~2.2]{OpenAI2026} in full.

For $m\geq1$, $u\in\R^m$, and $0<h<1$, put
\[
 X_h(u)=(h^{u_1},\ldots,h^{u_m}).
\]
A \emph{label} $\lambda=(r_\lambda,b_\lambda)$ consists of a slope
$r_\lambda\in\R^m$, independent of $h$, and an offset $b_\lambda(h)$
defined for all sufficiently small positive $h$. For a finite nonempty
family $\Lambda$ of labels, write
\[
 F_h(x)=\min_{\lambda\in\Lambda}
       \bigl(r_\lambda\cdot x+b_\lambda(h)\bigr).
\]
A label is \emph{active at $x$} if its affine value equals $F_h(x)$.
Every minimizing label is active when there is a tie. All offsets in a
finite family are understood to be restricted to a common interval on
which they are defined.

\begin{lemma}[Affine approximation with arbitrary tolerance]
\label{lem:affine}
Let $m\geq1$ be an integer, let $a<b$ be real numbers, and let
$E:[a,b]^m\to(0,\infty)$ be any function. There is a finite nonempty
family $\Lambda$ of labels with
\[
 r_\lambda\in[a,b]^m,
 \qquad b_\lambda(h)=o(h^a)\quad(h\to0+),
\]
such that the following implication holds. If $h_\nu\to0+$,
$p^\nu\in\R^m$ converges to $p\in[a,b]^m$, and a fixed label
$\lambda\in\Lambda$ is active at $X_{h_\nu}(p^\nu)$ for every $\nu$,
then
\[
 \norm{p-r_\lambda}_\infty<E(r_\lambda).
\]
The slopes and the family are independent of $h$. No regularity is
assumed of $E$, and the approximating vectors $p^\nu$ need not lie in
$[a,b]^m$.
\end{lemma}
